\begin{lemma}[The one-dimensional viscosity criterion for convexity]
Let $T>0$ and let $G:\R\to\R$ be continuous on $[0,T]$. Suppose that for every
$\ell_0\in(0,T)$ and all $A,B,C\in\R$, if the function
$\tau\mapsto G(\tau)-\bigl(A\tau^2+B\tau+C\bigr)$ has a local maximum on $[0,T]$ at
$\ell_0$, then $A\ge0$. Then $G$ is convex on $[0,T]$.
\end{lemma}
\begin{proof}
The interval $[0,T]$ is convex, so by the adjacent-slope characterisation of convexity it
suffices to show that for all $a,b\in[0,T]$ and all $\tau$ with $a<\tau<b$,
\[
  \frac{G(\tau)-G(a)}{\tau-a}\le\frac{G(b)-G(\tau)}{b-\tau}.
\]
Suppose this fails for some such $a,\tau,b$. Then $a<b$, and the three quantities
$b-a$, $\tau-a$, $b-\tau$ are all positive.

\emph{Step 1: the chord lies strictly below $G$ at $\tau$.}
Let
\[
  \lambda(r)=\frac{(b-r)G(a)+(r-a)G(b)}{b-a}
\]
be the affine function with $\lambda(a)=G(a)$ and $\lambda(b)=G(b)$. The failed
inequality, multiplied by the positive number $(\tau-a)(b-\tau)$, reads
\[
  \bigl(G(b)-G(\tau)\bigr)(\tau-a)<\bigl(G(\tau)-G(a)\bigr)(b-\tau),
\]
that is, after moving the terms containing $G(\tau)$ to one side,
\[
  G(b)(\tau-a)+G(a)(b-\tau)<G(\tau)\bigl((b-\tau)+(\tau-a)\bigr)=G(\tau)(b-a).
\]
Dividing by $b-a>0$ gives $\lambda(\tau)<G(\tau)$.

\emph{Step 2: the gap has a positive maximum.}
Put $h=G-\lambda$. Since $[a,b]\subseteq[0,T]$, by the hypothesis $G$ is continuous on
$[a,b]$, and $\lambda$ is continuous, so $h$ is continuous on the nonempty compact set
$[a,b]$ and attains a maximum there, say at $\tau_1$. Put $m=h(\tau_1)$. Since
$h(\tau)=G(\tau)-\lambda(\tau)>0$ by Step 1 and $h(\tau)\le m$, we get $m>0$.

\emph{Step 3: a strictly concave perturbation whose maximum is still positive.}
Put $\eta=m/(b-a)^2>0$ and
\[
  K(r)=h(r)+\eta\,(r-a)(r-b)\qquad(r\in[a,b]).
\]
For every $r\in[a,b]$ we have $(r-a)(r-b)\ge-(b-a)^2/4$: indeed
$(2r-a-b)^2\ge0$ expands to $4r^2-4r(a+b)+(a+b)^2\ge0$, while
$-4(r-a)(r-b)=-4r^2+4r(a+b)-4ab$, and adding the two gives
$-4(r-a)(r-b)\le(a+b)^2-4ab=(b-a)^2$. Multiplying by $\eta>0$ and using
$\eta(b-a)^2/4=m/4$,
\[
  K(\tau_1)=m+\eta(\tau_1-a)(\tau_1-b)\ge m-\frac m4=\frac{3m}{4}>\frac m2 .
\]
$K$ is continuous on $[a,b]$, so it attains a maximum there, say at $\tau_2$, and
$K(\tau_2)\ge K(\tau_1)>m/2>0$.

\emph{Step 4: the maximum is attained strictly inside.}
At the endpoints $\lambda$ agrees with $G$ and the perturbation vanishes, so
$K(a)=h(a)+0=G(a)-\lambda(a)=0$ and likewise $K(b)=0$. Since $K(\tau_2)>0$, we conclude
$\tau_2\ne a$ and $\tau_2\ne b$, so $a<\tau_2<b$. As $0\le a$ and $b\le T$ this gives
$\tau_2\in(0,T)$.

\emph{Step 5: the quadratic upper test.}
Expanding $\lambda$ and the perturbation, set
\[
  A=-\eta,\qquad
  B=\frac{G(b)-G(a)}{b-a}+\eta(a+b),\qquad
  C=\frac{bG(a)-aG(b)}{b-a}-\eta ab+K(\tau_2),
\]
so that for every $u\in\R$
\[
  A u^2+Bu+C=-\eta(u-a)(u-b)+\lambda(u)+K(\tau_2),
\]
using $\lambda(u)=\bigl[(G(b)-G(a))u+bG(a)-aG(b)\bigr]/(b-a)$. Consequently
\begin{equation}\label{eq:vcoi-shape}
  G(u)-\bigl(Au^2+Bu+C\bigr)
   =\bigl(G(u)-\lambda(u)\bigr)+\eta(u-a)(u-b)-K(\tau_2)
   =K(u)-K(\tau_2).
\end{equation}

The open interval $(a,b)$ is a neighbourhood of $\tau_2$, so for every $r$ in
$(a,b)\cap[0,T]$ we have $r\in[a,b]$ and hence $K(r)\le K(\tau_2)$ by the choice of
$\tau_2$; by \eqref{eq:vcoi-shape} this says
\[
  G(r)-\bigl(Ar^2+Br+C\bigr)\le0=G(\tau_2)-\bigl(A\tau_2^2+B\tau_2+C\bigr).
\]
Therefore $\tau\mapsto G(\tau)-\bigl(A\tau^2+B\tau+C\bigr)$ has a local maximum on
$[0,T]$ at $\tau_2$, and $\tau_2\in(0,T)$.

\emph{Step 6: the contradiction.}
The hypothesis applied with these $A,B,C$ and this $\ell_0=\tau_2$ gives $A\ge0$. But
$A=-\eta<0$, a contradiction. Hence the adjacent-slope inequality holds for all
admissible $a,\tau,b$, and $G$ is convex on $[0,T]$.
\end{proof}
